\chapter{分析引论}

按顺序记录每一天的用量，逐次计算一笔存款的本息，或将一次计算的结果作为下一次的输入，都会得到一列数。我们既关心第\(n\)个数是多少，也关心前\(n\)项合起来有多大；当项数不断增加时，还要研究它们是否接近某个确定的值。

本章先介绍数列的基本概念与等差、等比数列，再引入极限，研究有限步骤怎样确定无穷过程的性质。数列也是近似计算与数学分析的基本对象，后面的函数与微积分将继续使用这些思想。

\section{数列}
\label{sec:analysis-sequence-basics}

\subsection{数列的定义}

\begin{definition}
给每个正整数\(n\)指定唯一的实数\(a_n\)，按下标顺序排列得到
\begin{equation*}
 a_1,a_2,\ldots,a_n,\ldots,
\end{equation*}
称为一个{\heiti{}数列}，记作\(\{a_n\}\)。\(a_n\)是数列的第\(n\)项，\(n\)称为下标。若只取下标\(1,2,\ldots,N\)，便得到有限数列；本章讨论极限时，均指无穷数列。
\end{definition}

数列中的数可以重复，排列顺序也有意义。例如，\(1,-1,1,-1,\ldots\)与\(-1,1,-1,1,\ldots\)具有相同的取值集合，却是不同的数列。

若用关于\(n\)的式子直接给出\(a_n\)，称为数列的{\heiti{}通项公式}，例如\(a_n=2n-1\)。有了通项公式，就可以由下标确定相应的项。

记前\(n\)项的和为\(S_n\)，即
\begin{equation}\label{eq:analysis-partial-sum-definition}
 S_n=\sum_{k=1}^{n}a_k=a_1+a_2+\cdots+a_n.
\end{equation}
其中\(k\)只是求和下标，可以换成其他字母。约定\(S_0=0\)，则
\begin{equation}\label{eq:analysis-term-from-sum}
 a_n=S_n-S_{n-1},\qquad n\ge1.
\end{equation}
前\(n\)项的积类似地记作\(\prod_{k=1}^{n}a_k\)；特别地，\(n!=1\cdot2\cdots n\)，并约定\(0!=1\)。

\begin{example}
写出数列\(a_n=(-1)^{n-1}n\)的前五项。若另一数列的前\(n\)项和为\(S_n=n(n+1)\)，求它的通项。
\end{example}

只知道有限多项，通常不能唯一确定后面的项。例如，\(a_n=n\)与\(b_n=n+(n-1)(n-2)(n-3)\)的前三项相同，第四项却不同。观察可以帮助猜测通项，证明还须使用题目给出的完整条件。

\begin{example}
分别写出\(1,0,1,0,\ldots\)与\(1,-1/2,1/3,-1/4,\ldots\)的一个通项公式。对数列\(1,1/2,1/2,1/3,1/3,1/3,\ldots\)，其中\(1/k\)依次重复\(k\)次，试用下标的取值范围表示它的通项。
\end{example}

\subsection{单调性与有界性}

若对所有\(n\)都有\(a_{n+1}\ge a_n\)，称\(\{a_n\}\){\heiti{}单调递增}；将不等号换为\(\le\)，便是{\heiti{}单调递减}。取严格不等号时，称为严格单调。证明单调性通常考察差\(a_{n+1}-a_n\)的符号；对正数列，也可比较比值\(a_{n+1}/a_n\)与\(1\)的大小。

若有常数\(M\)，使所有项都满足\(a_n\le M\)，称数列有上界；下界类似。若上下界都有，称数列{\heiti{}有界}，也就是存在\(K>0\)，使\(|a_n|\le K\)对所有\(n\)成立。

有上界不表示上界一定是某一项。例如\(a_n=n/(n+1)\)严格递增，以\(1\)为上界，却没有最大项。数列是否单调、是否有界，将是后面判断收敛性的重要依据。

\begin{example}
判断\(a_n=(2n+1)/(n+1)\)、\(b_n=(-1)^n(1-1/n)\)、\(c_n=n(6-n)\)的单调性与有界性，并求它们存在的最大项、最小项。有界与单调能否互相推出？
\end{example}

\subsection{数学归纳法}

有些命题涉及每个正整数，逐个代入不能穷尽。若能从第一步出发，证明成立的一步总能推出下一步，就可以覆盖所有正整数。

\begin{theorem}[第一数学归纳法]
设\(P(n)\)是关于正整数\(n\)的命题。若\(P(1)\)成立，且对任意\(k\ge1\)，都有\(P(k)\Rightarrow P(k+1)\)，则\(P(n)\)对所有正整数成立。
\end{theorem}

归纳步骤暂时假设\(P(k)\)成立，是为了证明一个蕴含关系；它与已经验证的起始情形共同起作用。若命题从\(n_0\)开始，只须相应改变起点。

例如，\(1+3+\cdots+(2n-1)=n^2\)在\(n=1\)时成立。假设对\(n=k\)成立，再添下一项\(2k+1\)，便得\(k^2+2k+1=(k+1)^2\)，故公式对所有正整数成立。

\begin{example}
用数学归纳法证明\(\sum_{k=1}^n k(k+1)=n(n+1)(n+2)/3\)。再对\(x\ge-1\)证明Bernoulli不等式\((1+x)^n\ge1+nx\)，说明归纳步骤中为何可以将不等式两边乘以\(1+x\)。
\end{example}

有时，证明下一步需要使用前面若干步的结论，所用的下标也未必紧邻下一步。这时可以采用下面的形式。

\begin{theorem}[第二数学归纳法]
\label{thm:analysis-strong-induction}
设\(n_0\)为正整数，\(P(n)\)是关于整数\(n\ge n_0\)的命题。若\(P(n_0)\)成立，且对任意\(k\ge n_0\)，由\(P(n_0),P(n_0+1),\ldots,P(k)\)同时成立能够推出\(P(k+1)\)，则\(P(n)\)对所有\(n\ge n_0\)成立。
\end{theorem}

\begin{proof}
将“\(P(n_0),P(n_0+1),\ldots,P(k)\)均成立”记为\(Q(k)\)。由条件，\(Q(n_0)\)成立；若\(Q(k)\)成立，则\(P(k+1)\)成立，连同已有的结论便得到\(Q(k+1)\)。对\(Q(k)\)使用第一数学归纳法，即得结论。
\end{proof}

第二数学归纳法也称为{\heiti{}强归纳法}，归纳假设允许使用从起点到第\(k\)步的全部结论。两种形式可以互相推出，使用时依照证明中实际需要的结论选择即可。

\begin{example}
用第二数学归纳法证明：每个大于\(1\)的整数都可以写成若干个质数的乘积。
\end{example}

对于用到前两项的递推关系，可以先验证前两个起始情形，再同时使用这两项的结论作归纳推导；也可以将相邻两个命题合成一个命题，使用第一数学归纳法。起始情形需要验证几项，应当由递推所需的条件决定。

\begin{exercise}
\label{ex:analysis-induction}
用适当的数学归纳法完成下列各题。

（1）证明\(1^3+2^3+\cdots+n^3=[n(n+1)/2]^2\)。

（2）证明\(10^n-9n-1\)能被\(81\)整除。

（3）证明当\(n\ge4\)时，\(2^n\ge n^2\)。

（4）设\(a_1=1,a_2=2,a_{n+2}=a_{n+1}+a_n\)。证明\(n\le a_n\le2^{n-1}\)。

（5）设\(a_1=1\)，且当\(n\ge2\)时，\(a_n=\sum_{k=1}^{n-1}a_k\)。猜测并证明通项公式，再将给定关系化为相邻两项之间的递推关系。

（6）证明每个正整数都能表示为若干个互不相同的\(2\)的幂之和，其中指数为非负整数。
\end{exercise}


\newpage
\section{等差数列与等比数列}
\label{sec:analysis-arithmetic-geometric}

\subsection{等差数列}

\begin{definition}
若数列相邻两项的差\(a_{n+1}-a_n=d\)与\(n\)无关，称为{\heiti{}等差数列}，常数\(d\)称为公差。
\end{definition}

将\(a_2-a_1=d,\ldots,a_n-a_{n-1}=d\)相加，中间各项消去，得到
\begin{equation}\label{eq:analysis-arithmetic-term}
 a_n=a_1+(n-1)d.
\end{equation}
任意两项满足\(a_n-a_m=(n-m)d\)；下标和相同的两对项，其和也相同。将前\(n\)项正序、逆序各写一次，逐项相加，每一对都是\(a_1+a_n\)，所以
\begin{equation}\label{eq:analysis-arithmetic-sum}
 S_n=\frac{n(a_1+a_n)}2=na_1+\frac{n(n-1)}2d.
\end{equation}
特别地，\(1+2+\cdots+n=n(n+1)/2\)。等差数列的部分和是关于\(n\)的不超过二次的多项式，且常数项为零；反过来，由\(a_n=S_n-S_{n-1}\)也可以恢复等差数列。

\begin{example}
等差数列满足\(a_3+a_7=20,S_{10}=110\)。求通项与前\(n\)项和，再求使\(S_n>1000\)成立的最小正整数\(n\)。
\end{example}

\subsection{等比数列}

\begin{definition}
一个各项非零的数列，若相邻两项之比\(a_{n+1}/a_n=q\)与\(n\)无关，称为{\heiti{}等比数列}，常数\(q\)称为公比。
\end{definition}

将相邻两项的比连乘，得到
\begin{equation}\label{eq:analysis-geometric-term}
 a_n=a_1q^{n-1}.
\end{equation}
它描述固定倍率的变化；\(q<0\)时符号交替，因而固定倍率不一定产生单调数列。

设\(S_n=a_1+a_1q+\cdots+a_1q^{n-1}\)，乘以\(q\)后各项依次错开一位。两式相减，中间项消去，得\((1-q)S_n=a_1(1-q^n)\)，因而
\begin{equation}\label{eq:analysis-geometric-sum}
 S_n=a_1\frac{1-q^n}{1-q}\quad(q\ne1);
 \qquad S_n=na_1\quad(q=1).
\end{equation}
这种乘以公比使两组项对齐，再消去共同部分的方法，称为{\heiti{}错位相减}。

\begin{example}
等比数列满足\(a_2=6,a_5=162\)，求首项、公比及前\(n\)项和。另取\(b_n=(-1/2)^{n-1}\)，写出其部分和，比较原数列与部分和数列的变化。
\end{example}

\newpage
\section{极限}
\label{sec:analysis-limits}
数列\(a_n=n/(n+1)\)的每一项都小于\(1\)，但下标越大，它与\(1\)的差就越小。等比数列的求和公式又给出
\begin{equation*}
 \frac12+\frac1{2^2}+\cdots+\frac1{2^n}=1-\frac1{2^n}.
\end{equation*}
左边不断加入新的项，右边却表明，无论加到哪一项，所得的和都小于\(1\)，并且与\(1\)的差可以越来越小。求某一项或有限项的和，可以通过有限次运算完成；当下标不断增大时，我们还希望描述整个变化过程的性质。

这里的“接近”不只是看到小数点后的几位相同。若只要求误差小于\(0.01\)，可以取一批足够靠后的项；若要求小于\(0.0001\)，就需要把下标取得更大。真正要说明的是：无论给出怎样的正误差，都能找到一个下标，使它后面的所有项满足要求。极限将这种关系概括为一个确定的数学概念。

这个概念也不要求变化过程一定到达那个数。上面的两个数列都不取到\(1\)，但不妨碍我们用\(1\)描述它们的最终趋势；数列\((-1)^n/n\)在零的两侧交替变化，也可以越来越接近零。因此，要控制的是与目标值的距离，而不是每一步变化的方向。

函数也有类似的问题。考察\(x\)接近\(1\)时\(x^2\)的变化，我们关心的是：怎样限制\(x\)与\(1\)的距离，才能使\(x^2\)与\(1\)的误差达到要求？数列中取足够大的下标，函数中取足够小的距离，都是用自变量的范围控制结果的误差。

有了这样的表述，有限的估计就能够说明无限过程的性质。计算若干项可以帮助我们观察、猜测，极限的定义则要给出对全部后续项有效的保证。下面将“任意给定的精度”和“保证这个精度的范围”分别用数学符号表示。

\subsection{极限的定义}
\begin{definition}[数列的极限]
设\(\{a_n\}\)是一个数列，\(a\)为实数。若
\begin{equation}\label{eq:analysis-sequence-limit-definition}
 \forall\varepsilon>0,\ \exists N\in\mathbb{N}^{*},\quad
 n>N\ \Longrightarrow\ |a_n-a|<\varepsilon,
\end{equation}
则称\(\{a_n\}\)的极限为\(a\)，或称\(\{a_n\}\){\heiti{}收敛于}\(a\)，记作\(\lim_{n\to+\infty}a_n=a\)。
\end{definition}

这里的\(N\)可以依赖于\(\varepsilon\)；一旦选定\(N\)，它后面的所有项都应满足误差要求。这种定义称为{\heiti{}\(\varepsilon\)-\(N\)定义}。对函数，控制误差的条件改为自变量与指定点足够接近。

\begin{definition}[函数的极限]
设函数\(f\)在\(x_0\)的去心邻域内有定义，\(L\)为实数。若
\begin{equation}\label{eq:analysis-function-limit-definition}
 \forall\varepsilon>0,\ \exists\delta>0,\quad
 0<|x-x_0|<\delta\ \Longrightarrow\ |f(x)-L|<\varepsilon,
\end{equation}
则称\(f(x)\)在\(x\to x_0\)时以\(L\)为极限，记作\(\lim_{x\to x_0}f(x)=L\)。
\end{definition}

这称为{\heiti{}\(\varepsilon\)-\(\delta\)定义}。条件\(0<|x-x_0|<\delta\)只考察点附近的值，不要求\(f(x_0)\)有定义。数列中取足够大的下标，函数中取足够小的距离，都是对任意给定的精度，找到能够保证它的范围。

\subsection{无穷大量与无穷小量}
无穷大量与无穷小量描述的是量在指定极限过程中的变化。

\begin{definition}[无穷大量]
设函数\(f\)在\(x_0\)的去心邻域内有定义。若
\begin{equation}
 \forall M>0,\ \exists\delta>0,\quad
 0<|x-x_0|<\delta\ \Longrightarrow\ |f(x)|>M,
\end{equation}
则称\(f(x)\)是当\(x\to x_0\)时的{\heiti{}无穷大量}，记作\(\lim_{x\to x_0}|f(x)|=+\infty\)。
\end{definition}

若将定义中的\(|f(x)|>M\)换为\(f(x)>M\)，则称\(f(x)\)趋于\(+\infty\)；换为\(f(x)<-M\)，则称\(f(x)\)趋于\(-\infty\)。例如，当\(x\to0\)时，\(1/x^2\)趋于\(+\infty\)，\(1/x\)的绝对值趋于无穷，但它本身不趋于\(+\infty\)或\(-\infty\)。

\begin{definition}[无穷小量]
在指定的极限过程中趋于\(0\)的量称为{\heiti{}无穷小量}。例如，若\(\lim_{x\to x_0}f(x)=0\)，则称\(f(x)\)是当\(x\to x_0\)时的无穷小量。
\end{definition}

无穷小量的误差条件就是极限定义中目标值为\(0\)的情形。将上述条件中的\(0<|x-x_0|<\delta\)换为\(n>N\)，便得到数列的相应定义；例如，\(1/n\)是当\(n\to+\infty\)时的无穷小量。

数列趋于\(+\infty\)或\(-\infty\)时，称它有相应的广义极限。仅有\(|a_n|\to+\infty\)不保证符号最终固定，例如\(a_n=(-1)^n n\)。下面将有限极限与有确定符号的广义极限区分使用。

\newpage
\section{数列极限}
\label{sec:analysis-sequence-limits}
上面已经介绍了数列极限的基本定义，其他形式的等价定义请参考教材叙述。

有极限的数列称为{\heiti{}收敛数列}，不收敛的数列称为{\heiti{}发散数列}。

\begin{definition}[子列]
设\(\{ a_{n} \}\)是一个数列，\(\{ n_{k} \}\)是一个\textbf{严格单调递增}的自然数列（即\(\forall k \in \mathbb{N}^{*} , n_{k} \in \mathbb{N}^{*}\)），则可以得到另一个数列：
\begin{equation*}
    a_{n_{1}} , a_{n_{2}} , \ldots , a_{n_{k}} , \cdots
\end{equation*}
记为\(\left\{ a_{n_{k}} \right\}\)，称\(\left\{ a_{n_{k}} \right\}\)为\(\left\{ a_{n} \right\}\)的一个子列。
\end{definition}

显然，“子列”一词说明，子列中的元素均来源于“父列”，即原数列；且子列不改变元素在原数列中的排列顺序。来源于同一数列的子列可以有无穷多个，数列自身也是该数列的一个子列。（此处可参考集合里的“子集”概念）

\subsection{放缩法}
放缩法是用不等式控制所研究的量，或控制它与目标值之间的误差。例如，要证明\(a_n\to a\)，可以寻找容易估计的\(b_n\)，使\(|a_n-a|\le b_n\)，再对任意给定的\(\varepsilon>0\)，确定使\(b_n<\varepsilon\)成立的下标范围。

放缩的方向与精度应当由目标决定。估计式既要容易处理，也要保留足够的信息，使所需误差能够趋于零。对同一个表达式，证明有界与证明收敛，所需的估计可能不同。

第五章将借助导数、中值定理与凸性，介绍构造上下界与控制误差的方法；选讲中的Taylor公式还会进一步给出近似的精度。这里先从具体表达式出发，练习怎样得到有效的估计。

\begin{example}
    以下证明请使用极限的定义进行证明，不要使用其他定理或推论。

    （1）数列\(\{n^5/2^n\}\)收敛于\(0\)。

    （2）数列\(\{\sqrt[n]{n}\}\)的极限是\(1\)。

    （3）\(\displaystyle\lim_{n\to\infty}\frac{3n^2}{n^2-4}=3\)。

    （4）\(\displaystyle\lim_{n\to\infty}\frac{a^n}{n!}=0\quad(a>0)\)。
\end{example}

\begin{example}
利用有理化建立误差上界，用定义证明\(\sqrt{n+1}-\sqrt n\to0\)，并对给定的\(\varepsilon>0\)写出一个可用的\(N\)。
\end{example}
\newpage
\subsection{收敛数列的基本性质}
\begin{theorem}\label{thm:analysis-sequence-limit-properties}
设\(a_n\to a,b_n\to b\)，则有以下性质。

（1）极限唯一；改变有限多项，不改变收敛性及极限。

（2）收敛数列有界；它的每个子列都收敛于同一极限。

（3）若从某项起\(a_n\le b_n\)，则\(a\le b\)。若\(a>0\)，则从某项起\(a_n>0\)；负极限类似。

（4）对常数\(\alpha,\beta\)，有
\begin{equation*}
 \alpha a_n+\beta b_n\to\alpha a+\beta b,\qquad
 a_nb_n\to ab,\qquad
 \frac{a_n}{b_n}\to\frac ab\quad(b\ne0).
\end{equation*}
\end{theorem}

\begin{proof}
若同一数列有不同极限\(a,b\)，取\(\varepsilon=|a-b|/3\)，当下标充分大时，两种误差估计与三角不等式矛盾。改变有限多项时，只须把\(N\)取得超过这些项的下标，原来的误差估计仍然成立。取\(\varepsilon=1\)，充分靠后的项都在\((a-1,a+1)\)内，再将有限多个前项计入，就得到全数列的界。子列下标严格增加，因而最终也进入原数列保证误差的范围。

加法与数乘用三角不等式控制误差；乘法则分解为
\begin{equation*}
 a_nb_n-ab=a_n(b_n-b)+b(a_n-a),
\end{equation*}
再使用\(a_n\)有界。商的情形，由\(b_n\to b\ne0\)，充分靠后的项满足\(|b_n|>|b|/2\)，于是
\begin{equation*}
 \left|\frac{a_n}{b_n}-\frac ab\right|
 \le\frac{2}{|b|}|a_n-a|+\frac{2|a|}{|b|^2}|b_n-b|.
\end{equation*}
右边两项都能任意小，因而商也趋于相应的极限。

若\(a>0\)，取\(\varepsilon=a/2\)，便有充分靠后的项满足\(a_n>a/2>0\)。负极限类似。若从某项起\(a_n\le b_n\)，而\(a>b\)，则\(a_n-b_n\to a-b>0\)，由保号性得\(a_n>b_n\)最终成立，产生矛盾，故\(a\le b\)。
\end{proof}

对于非负数列，若\(a_n\to a\)，则\(\sqrt{a_n}\to\sqrt a\)。可由\(|\sqrt{a_n}-\sqrt a|\le\sqrt{|a_n-a|}\)直接控制误差。后面在含平方根的递推式中取极限时将使用这一点。

{\kaiti{}思考\quad （1）设\(\left\{ a_{n} \right\}\)收敛，\(\left\{ b_{n} \right\}\)发散，问：数列\(\left\{ a_{n} + b_{n} \right\} , \left\{ a_{n} b_{n} \right\}\)的敛散性如何？

（2）设\(\left\{ a_{n} \right\} , \left\{ b_{n} \right\}\)均发散，问：数列\(\left\{ a_{n} + b_{n} \right\} , \left\{ a_{n} b_{n} \right\}\)的敛散性如何？

（3）设\(a_{n} < b_{n} < c_{n} , n \in \mathbb{N}^{*}\)，若\(\lim_{n \to \infty} \left( c_{n} - a_{n} \right) = 0\)，问：数列\(b_{n}\)是否收敛？

（4）找出下列运算中的错误：
\begin{equation*}
    \lim_{n \to \infty} \left( \frac{1}{n + 1} + \cdots + \frac{1}{2 n} \right) = \lim_{n \to \infty} \left( \frac{1}{n + 1} \right) + \cdots + \lim_{n \to \infty} \left( \frac{1}{2 n} \right) = 0 + \cdots + 0 = 0
\end{equation*}

（5）设已知\(\left\{ a_{n} \right\}\)收敛于\(a\)，又对每一个\(n\)有\(b < a_{n} < c\)，问：是否成立\(b < a < c\)？

（6）设已知\(\left\{ a_{n} \right\}\)收敛于\(a\)，又有\(b \le a \le c\)，问：是否有\(\exists N \in \mathbb{N}^{*} , \text{s.t.} \forall n > N , b \le a_{n} \le c\)成立？

（7）设已知\(\lim_{n \to \infty} a_{n} = 0\)，问：是否有\(\lim_{n \to \infty} \left( a_{1} a_{2} \cdots a_{n} \right) = 0\)？反之如何？}

\begin{example}

    利用收敛数列的性质完成以下证明或计算：

    （1）证明：若一数列收敛，则该数列必有最大值或最小值，但不一定同时有最大值和最小值。

    （2）若\(\lim_{n \to \infty} \frac{x_{n} - a}{x_{n} + a} =
    0\)，证明\(\lim_{n \to \infty} x_{n} = a\).

    （3）设\(s_{n} = a + 3 a^{2} + \cdots + \left( 2 n - 1 \right) a^{n} , \left| a
    \right| < 1\)，求\(\left\{ s_{n} \right\}\)的极限。

    （4）求\(\lim_{n \to \infty} \prod _{i = 0}^{n} \left( 1 + x^{2^{i}}
    \right)\)，其中\(\left| x \right| < 1\).
\end{example}

\subsection{数列收敛性的研究工具}
\paragraph{1°夹逼原理} \mbox{}

难以直接处理一个数列时，可以寻找两个较为简单的数列，从上下两边限制它的取值。

\begin{theorem}[夹逼原理]
\label{thm:analysis-squeeze}
设从某项起有\(a_n\le b_n\le c_n\)，且\(a_n\to a,c_n\to a\)，则\(b_n\to a\)。
\end{theorem}

\begin{proof}
对任意\(\varepsilon>0\)，将\(N\)取得充分大，则当\(n>N\)时，夹逼关系成立，且\(|a_n-a|<\varepsilon,|c_n-a|<\varepsilon\)。于是
\begin{equation*}
 a-\varepsilon<a_n\le b_n\le c_n<a+\varepsilon,
\end{equation*}
即\(|b_n-a|<\varepsilon\)，故\(b_n\to a\)。
\end{proof}

特别地，若\(|b_n-a|\le d_n\)，且\(d_n\to0\)，对\(0\le|b_n-a|\le d_n\)使用夹逼原理，就得到\(b_n\to a\)。前面的误差估计由此也可以直接用于判断收敛。

\begin{example}
    设\(a > 0 , b > 0\)，求极限\(\lim_{n \to \infty} {\left( a^{n} + b^{n} \right)}^{\frac{1}{n}}\).
\end{example}

\newpage
\paragraph{2°单调有界原理} \mbox{}

\begin{theorem}[单调有界原理]
\label{thm:analysis-monotone-convergence}
单调递增且有上界的实数列收敛；单调递减且有下界的实数列也收敛。
\end{theorem}

本章将这一结论作为实数的基本性质使用。它保证满足条件的数列有极限，即使暂时还不知道极限是多少，也可以先确定收敛性，再研究极限的值。

若单调递增数列无上界，则对任意\(M>0\)，总能找到一项大于\(M\)，此后各项均不小于该项，故数列趋于\(+\infty\)。同理，单调递减而无下界的数列趋于\(-\infty\)。

\begin{example}
    证明：若\(e_{n} = {\left( 1 + \frac{1}{n} \right)}^{n}\)，\(n \ge 1\)，则数列\(\left\{ e_{n} \right\}\)收敛。
\end{example}

\paragraph{3°区间套定理} \mbox{}

单调有界原理也能把逐步缩小的区间确定为一个点。设闭区间逐层包含，且长度趋于零，便有如下结论。

\begin{theorem}[闭区间套定理]\label{thm:analysis-nested-intervals}
设\([a_{n+1},b_{n+1}]\subseteq[a_n,b_n]\)，且\(b_n-a_n\to0\)。则存在唯一实数\(\xi\)，属于全部区间，并有\(a_n\to\xi,b_n\to\xi\)。
\end{theorem}

\begin{proof}
\(a_n\)递增、\(b_n\)递减，且\(a_1\le a_n\le b_n\le b_1\)，故分别收敛。两极限之差为零，记共同极限为\(\xi\)。对任意固定的\(n\)，后面各项仍在\([a_n,b_n]\)中，取极限得\(a_n\le\xi\le b_n\)。若另有公共点，它与\(\xi\)的距离不超过每个区间的长度，故只能相同。
\end{proof}

只要保留条件能持续使用，反复二分就得到长度依次减半的区间套。若初始区间长度为\(L\)，二分\(k\)次后用中点\(x_k\)近似公共点，则误差不超过\(L/2^{k+1}\)。给定精度\(\varepsilon>0\)，由\(L/2^{k+1}<\varepsilon\)就能预先确定二分次数。这既是存在性证明，也是带有误差保证的计算方法。

\begin{theorem}\label{thm:analysis-bounded-subsequence}
每个有界实数列都有一个收敛子列。
\end{theorem}

\begin{proof}
取一个包含全部项的闭区间。把它二分，至少一半包含无穷多个下标对应的项；保留这一半并继续二分，就得到长度趋于零的区间套。由定理\ref{thm:analysis-nested-intervals}，其公共点唯一。每次从保留区间中选一项，并使下标大于上一次，所选项与公共点的距离不超过该区间长度，因而构成收敛子列。
\end{proof}

这里得到的只是一个收敛子列，并不保证原数列收敛。它与前面的“收敛数列的所有子列同极限”有不同的作用。第二章将用这些方法研究连续函数在闭区间上的性质。

\newpage
\paragraph{4°Cauchy收敛准则} \mbox{}

Cauchy收敛准则是判断数列是否收敛/发散的充要条件。

\begin{definition}[基本列（柯西列）]
若对任意\(\varepsilon>0\)，存在\(N\)，使任意\(m,n>N\)都有\(|a_m-a_n|<\varepsilon\)，则称\(\{a_n\}\)为{\heiti{}基本列}。
\end{definition}
\begin{samepage}
\begin{theorem}[Cauchy收敛准则]
    数列\(\left\{ a_{n} \right\}\)收敛的充要条件是\(\left\{ a_{n} \right\}\)是基本列。
\end{theorem}

\begin{proof}
收敛数列的两项在下标充分大时都接近极限，三角不等式即给出必要性。反过来，基本列取误差\(1\)便知其尾部有界，因而全数列有界。由定理\ref{thm:analysis-bounded-subsequence}，可取子列\(a_{n_k}\to a\)。对给定\(\varepsilon>0\)，先使所有足够靠后的两项之差小于\(\varepsilon/2\)，再选其中一个子列项，使它与\(a\)的差也小于\(\varepsilon/2\)。于是所有足够靠后的项都满足\(|a_n-a|<\varepsilon\)，故原数列收敛。
\end{proof}
\end{samepage}

\begin{example}
    （1）设\(a_{n} = \sum _{i = 1}^{n} \frac{\sin i}{i^{2}}\)，证明：\(\left\{ a_{n} \right\}\)收敛。

    （2）证明：单调数列\(\left\{ a_{n} \right\}\)收敛的充要条件是它有一个收敛子列。
\end{example}

\subsection{判断数列发散的方法}
判断发散可以考察以下情形：

（1）数列无界；（2）有一个发散子列；（3）有两个子列收敛于不同的极限。

数列收敛的定义及Cauchy收敛准则仍然可以用来判断数列是否发散。

\par\noindent\begin{minipage}{\linewidth}
\setlength{\parindent}{2em}\setlength{\parskip}{6pt}
\begin{example}
    完成以下证明：

    （1）设\(S_{n} = 1 + \frac{1}{2} + \frac{1}{3} + \cdots + \frac{1}{n} , n \in
    \mathbb{N}^{*}\)，证明数列\(\left\{ S_{n} \right\}\)发散。

    （2）证明数列\(\left\{ \sin n \right\}\)发散。
\end{example}
\end{minipage}

\begin{exercise}
\label{ex:analysis-convergence-tools}
\mbox{}\par\nobreak

（1）用夹逼原理求\(\displaystyle\lim_{n\to\infty}\sum_{k=1}^{n}\frac1{\sqrt{n^2+k}}\)。

（2）设\(a_n=\sum_{k=1}^{n}(-1)^{k-1}/k\)。考察奇、偶子列的单调性，证明\(\{a_n\}\)收敛。

（3）若\(|a_{n+1}-a_n|\le2^{-n}\)，证明\(\{a_n\}\)收敛。若只知\(a_{n+1}-a_n\to0\)，能否得到同一结论？
\end{exercise}

\newpage
\section{Cauchy命题与Stolz定理}
\label{sec:analysis-cesaro-stolz}

\subsection{Cauchy命题}
\begin{theorem}
    设\(\{ x_{n} \}\)收敛于\(l\)，则它的前\(n\)项的平均值（所形成的数列）也收敛于\(l\)，即
\begin{equation}
    \lim_{n \to + \infty} \frac{\sum _{i = 1}^{n} x_{i}}{n} = l
\end{equation}
\end{theorem}
\paragraph{证明思路} \mbox{}

夹逼定理、单调有界原理等似乎都不好证明该命题，因此需要找些新的路径。首先观察式子
\begin{equation}
    \frac{x_{1} + x_{2} + \cdots + x_{n}}{n}
\end{equation}
如果每一项都与\(l\)充分接近，则其算术平均数也会与\(l\)接近。而\(\lim_{n \to \infty} x_{n} = l\)，只要令\(\varepsilon\)充分小，则从某项开始，后面的所有项将与\(l\)充分接近，即满足\(\left| x_{n} - l \right| < \varepsilon\)。因此对原式进行拆分：
\begin{equation}
    \frac{x_{1} + x_{2} + \cdots + x_{n}}{n} = \frac{x_{1} + \cdots + x_{N}}{n} + \frac{x_{N + 1} + \cdots + x_{n}}{n}
\end{equation}
而（为看清思路，先设各项\(x_{n} > 0\)，一般情形见下面的证明）
\begin{equation}
    \frac{x_{1} + \cdots + x_{N}}{n} \le \max \left\{ x_{1} , \ldots , x_{N} \right\} \times \frac{N}{n}
\end{equation}
\begin{equation}
    \left| \frac{x_{N + 1} + \cdots + x_{n}}{n - N} - l \right| \le \varepsilon
\end{equation}
因此，只要\(n\)取得足够大：前\(N\)项被\(n\)除以后可以任意小；后\(n - N\)项的平均值与\(l\)的差不超过\(\varepsilon\)，且它在整个算术平均值中的权重\(\frac{n - N}{n}\)任意接近1。也就是说：先由\(\lim_{n \to \infty} x_{n} = l\)取定\(N\)，再把\(n\)取得充分大，使\(\frac{N}{n}\)任意小。

从而确定两步走的思路，即先确定\(N\)，根据\(N\)对数列进行分拆，从而确定\(n\)，使得该式成立。

\begin{proof}

根据条件\(\lim_{n \to \infty} x_{n} = l\)，可知可以对任意给定的\(\varepsilon > 0\)，取定\(N\)，使得当\(n > N\)时，\(\left| x_{n} - l \right| < \varepsilon / 2\)成立。于是当\(n > N\)时，有如下估计：
\begin{multline}
    \left| \frac{x_{1} + x_{2} + \cdots + x_{n}}{n} - l \right| = \left| \frac{\left( x_{1} - l \right) + \left( x_{2} - l \right) + \cdots + \left( x_{n} - l \right)}{n} \right| \\
    \le \frac{\left| \left( x_{1} - l \right) + \cdots + \left( x_{N} - l \right) \right|}{n} + \frac{\left| x_{N + 1} - l \right| + \cdots + \left| x_{n} - l \right|}{n} < \frac{M}{n} + \frac{n - N}{n} \cdot \varepsilon / 2
\end{multline}

这里\(M = \left| \left( x_{1} - l \right) + \cdots + \left( x_{N} - l \right) \right|\)是一个确定的数，由此，只要取得
\begin{equation}
    N_{1} = \max \left\{ N , \left[ \frac{M}{\varepsilon / 2} \right] \right\}
\end{equation}

就能保证当\(n > N_{1}\)时，下式成立：
\begin{equation}
    \left| \frac{x_{1} + x_{2} + \cdots + x_{n}}{n} - l \right| < \varepsilon
\end{equation}

\end{proof}

\begin{example}
    （1）设\(\lim_{n \to \infty} x_{n} = + \infty\)，证明：\(\lim_{n \to \infty} \frac{\sum _{i = 1}^{n} x_{i}}{n} = + \infty\).

    （2）设\(\left\{ a_{n} \right\}\)单调递增，\(\lim_{n \to \infty} \frac{\sum _{i = 1}^{n} a_{i}}{n} = a\)，证明：\(\left\{ a_{n} \right\}\)收敛到\(a\).

    （3）已知\(k \in \mathbb{N}^{*}\)，设\(\left\{ a_{2 k - 1} \right\}\)收敛到\(a\)，\(\left\{ a_{2 k} \right\}\)收敛到\(b\)，且\(a \ne b\)，求\(\lim_{n \to \infty} \frac{\sum _{i = 1}^{n} a_{i}}{n}\).

    （4）设\(\left\{ a_{n} \right\}\)为正数列，且收敛到\(A\)，证明：\(\lim_{n \to \infty} {\left( \prod _{i = 1}^{n} a_{i} \right)}^{\frac{1}{n}} = A\).
\end{example}

\subsection{Stolz定理}
\paragraph{1°\texorpdfstring{\(0/0\)}{0/0}型} \mbox{}

设\(\left\{ a_{n} \right\} , \left\{ b_{n} \right\}\)都是无穷小量，其中\(\left\{ a_{n} \right\}\)是严格单调递减数列，又存在
\begin{equation}
    \lim_{n \to \infty} \frac{b_{n + 1} - b_{n}}{a_{n + 1} - a_{n}} = l
\end{equation}
则有
\begin{equation}
    \lim_{n \to \infty} \frac{b_{n}}{a_{n}} = l
\end{equation}
其中\(l\)为有限值或\(\pm \infty\).

\paragraph{2°\texorpdfstring{\(* / \infty\)}{*/∞}型} \mbox{}

设\(\left\{ a_{n} \right\}\)是严格单调递增的无穷大量，又存在
\begin{equation}
    \lim_{n \to \infty} \frac{b_{n + 1} - b_{n}}{a_{n + 1} - a_{n}} = l
\end{equation}
则有
\begin{equation}
    \lim_{n \to \infty} \frac{b_{n}}{a_{n}} = l
\end{equation}

其中\(l\)为有限值或\(\pm \infty\).

这两个定理都用相邻增量的比控制原量的比。以第二个定理的有限极限为例，当\(k\)充分大时，
\begin{equation*}
 (l-\varepsilon)(a_{k+1}-a_k)
 \le b_{k+1}-b_k
 \le(l+\varepsilon)(a_{k+1}-a_k).
\end{equation*}
从固定的\(N\)累加到\(n-1\)，再除以\(a_n\)，有限初始部分的影响趋于零，便把\(b_n/a_n\)夹在任意接近\(l\)的范围内。对于\(0/0\)型，改用正的增量\(a_k-a_{k+1}\)，从\(n\)累加到\(m\)，再令\(m\to\infty\)，利用\(a_m,b_m\to0\)即可。确定符号的无穷极限只需相应的单边估计。

    {\kaiti{}思考\quad 若将以上三个命题中的极限值\(l\)改为不带符号的无穷大\(\infty\)，则结论均不成立。试举出反例。

\enlargethispage{2\baselineskip}}

\enlargethispage{\baselineskip}
\begin{example}
    （1）设\(a_{n} = \frac{\sum _{i = 1}^{n} i !}{n !} , n \in \mathbb{N}^{*}\)，求\(\lim_{n
        \to \infty} a_{n}\).

    （2）已知\(\lim_{n \to \infty} a_{n} = a\)，证明：\(\lim_{n \to
        \infty} \frac{1}{2^{n}} \sum _{k = 0}^{n} \left( \begin{matrix}n \\ k\end{matrix} \right) a_{k} = a\).

    （3）设\(a_{1} > 0 , a_{n + 1} = a_{n} + \frac{1}{a_{n}} , n \in
    \mathbb{N}^{*}\)，证明：\(\lim_{n \to \infty} \frac{a_{n}}{\sqrt{2 n}} = 1\).

    （4）若\(\lim_{n \to \infty} a_{n} = a , \lim_{n \to \infty}
    b_{n} = b\)，证明：\(\lim_{n \to \infty} \frac{\sum _{i = 1}^{n} a_{i}
        b_{n - i + 1}}{n} = a b\).

    （5）设\(\left\{ a_{n} \right\}\)为正数列，且有\(\lim_{n \to \infty} \frac{a_{n
                + 1}}{a_{n}} = l\)，证明：\(\lim_{n \to \infty} \sqrt[n]{a_{n}} = l\).
    反之如何？
\end{example}


\newpage
\section{数项级数}
\label{sec:analysis-series}

有限个数可以逐项相加。若不断加入新的项，所得结果是否接近一个确定的数，就要研究部分和数列的极限。数项级数以数列为基础，将求和问题与收敛问题联系起来。

\subsection{数项级数的概念}

\begin{definition}
给定实数列\(\{a_n\}\)，表达式
\begin{equation*}
 \sum_{n=1}^{\infty}a_n=a_1+a_2+\cdots+a_n+\cdots
\end{equation*}
称为以\(a_n\)为通项的{\heiti{}数项级数}。记\(S_n=\sum_{k=1}^n a_k\)，称\(S_n\)为级数的部分和。若部分和数列收敛于\(S\)，则称级数{\heiti{}收敛}，\(S\)称为它的和，记作\(\sum_{n=1}^{\infty}a_n=S\)；否则称级数发散。
\end{definition}

这里的无穷求和由有限和的极限定义。不能先把它当作已经完成的有限次加法，再任意使用有限和的运算规则。

例如，等比级数在\(|q|<1\)时有
\begin{equation}\label{eq:analysis-geometric-series}
 \sum_{n=1}^{\infty}a_1q^{n-1}
 =\lim_{n\to\infty}a_1\frac{1-q^n}{1-q}
 =\frac{a_1}{1-q}.
\end{equation}
取\(a_1=1,q=1/2\)，便有\(1+1/2+1/4+\cdots=2\)。部分和永远小于\(2\)，却可以任意接近\(2\)，两者并不矛盾。

\begin{proposition}
若\(\sum a_n\)收敛，则\(a_n\to0\)。改变级数的有限项不改变其敛散性。若\(\sum a_n=A\)、\(\sum b_n=B\)，则对常数\(\alpha,\beta\)，有
\begin{equation*}
 \sum_{n=1}^{\infty}(\alpha a_n+\beta b_n)=\alpha A+\beta B.
\end{equation*}
\end{proposition}

\begin{proof}
由\(a_n=S_n-S_{n-1}\)知，收敛时\(a_n\to S-S=0\)。改变有限项，只使后面的部分和增加一个固定常数。最后一个结论由有限和的线性性质与数列极限的运算法则得到。
\end{proof}

通项趋于零只是必要条件。例如调和级数\(\sum 1/n\)发散，前面已通过其部分和说明这一点。通项本身的极限与部分和的极限，是两个不同的问题。

由部分和数列的Cauchy准则，级数收敛等价于：对任意\(\varepsilon>0\)，存在\(N\)，使任意\(m>n>N\)都有
\begin{equation}\label{eq:analysis-series-cauchy}
 \left|\sum_{k=n+1}^{m}a_k\right|<\varepsilon.
\end{equation}

\begin{example}
讨论\(\sum_{n=1}^{\infty}q^{n-1}\)在实数\(q\)取不同值时的敛散性。再考察\(1-1+1-1+\cdots\)，比较它的部分和与将相邻两项先合并后得到的部分和。
\end{example}

\subsection{从有限求和到无穷求和}

\paragraph{1°累加与裂项} \mbox{}

有限和可以分组、换序、提出常数。例如
\begin{equation*}
 \sum_{k=1}^n(\alpha a_k+\beta b_k)
 =\alpha\sum_{k=1}^n a_k+\beta\sum_{k=1}^n b_k.
\end{equation*}
若一项可以写成两个相邻量之差，累加后便只留下首尾：
\begin{equation}\label{eq:analysis-telescoping-sum}
 \sum_{k=1}^n(b_k-b_{k+1})=b_1-b_{n+1}.
\end{equation}
例如，\(1/[k(k+1)]=1/k-1/(k+1)\)，所以
\begin{equation*}
 \sum_{k=1}^n\frac1{k(k+1)}=1-\frac1{n+1},
 \qquad \sum_{k=1}^{\infty}\frac1{k(k+1)}=1.
\end{equation*}
这种方法称为{\heiti{}裂项相消}。它先处理每个有限部分和，再对部分和取极限。

差分也可以用来求有限和。将\((k+1)^3-k^3=3k^2+3k+1\)从\(k=1\)到\(n\)相加，利用自然数求和公式，得到
\begin{equation*}
 \sum_{k=1}^n k^2=\frac{n(n+1)(2n+1)}6.
\end{equation*}
求和公式因此可以来自代数恒等式，而不必逐条记忆。

\paragraph{2°错位相减} \mbox{}

等比数列的求和依靠错位相减。若每项还带有下标的权重，错位后权重的次数会降低。设\(T_n=\sum_{k=1}^n kq^{k-1}\)，其中\(q\ne0,1\)，则
\begin{equation*}
 (1-q)T_n=1+q+\cdots+q^{n-1}-nq^n.
\end{equation*}
由等比数列求和公式，得
\begin{equation}\label{eq:analysis-weighted-geometric-sum}
 T_n=\frac{1-(n+1)q^n+nq^{n+1}}{(1-q)^2}.
\end{equation}
当\(|q|<1\)时，\(nq^n\to0\)，于是
\begin{equation*}
 \sum_{k=1}^{\infty}kq^{k-1}=\frac1{(1-q)^2}.
\end{equation*}
若权重是更高次的多项式，仍可利用前后权重之差降低次数。

\begin{example}
求下列级数的和，并写出前\(n\)项和与其极限之间的差：
\begin{equation*}
 \text{（1）}\ \sum_{k=1}^{\infty}\frac1{k(k+2)};
 \qquad
 \text{（2）}\ \sum_{k=1}^{\infty}\frac{k}{2^k}.
\end{equation*}
\end{example}

\subsection{正项级数与比较判别}

若各项\(a_n\ge0\)，称\(\sum a_n\)为{\heiti{}正项级数}。其部分和单调不减，故由单调有界原理，正项级数收敛的充要条件是部分和有上界。这里首先需要控制总和，不一定先求出它的值。

\begin{theorem}[比较判别法]
设从某项起\(0\le a_n\le b_n\)。若\(\sum b_n\)收敛，则\(\sum a_n\)收敛；若\(\sum a_n\)发散，则\(\sum b_n\)发散。
\end{theorem}

\begin{proof}
略去有限项后，可设比较关系对所有项成立。若\(\sum b_n\)收敛于\(B\)，则\(\sum_{k=1}^n a_k\le\sum_{k=1}^n b_k\le B\)，故前者部分和有上界。第二个结论是第一个结论的逆否命题。
\end{proof}

例如对\(n\ge2\)，有\(1/n^2\le1/[n(n-1)]\)，而右侧级数可以裂项求和，故\(\sum 1/n^2\)收敛。这里并没有求出它的和。

若\(a_n,b_n>0\)，且\(a_n/b_n\to c\)，其中\(0<c<\infty\)，则充分靠后的项满足\((c/2)b_n\le a_n\le2cb_n\)，故两个级数同敛散。这称为比较判别法的极限形式。

\paragraph{\texorpdfstring{\(p\)}{p}级数} \mbox{}

级数\(\sum_{n=1}^{\infty}1/n^p\)在\(p>1\)时收敛，在\(p\le1\)时发散。\(p\le1\)时，\(1/n^p\ge1/n\)，可与调和级数比较。对\(p>1\)，将下标按\(2^j\le n<2^{j+1}\)分组，每组有\(2^j\)项，且
\begin{equation*}
 \sum_{n=2^j}^{2^{j+1}-1}\frac1{n^p}
 \le 2^j\frac1{2^{jp}}=\left(2^{1-p}\right)^j.
\end{equation*}
右边构成公比小于\(1\)的等比级数，便给出了全部部分和的上界。这是把逐项比较改为分组比较。

\begin{example}
判断下列级数的敛散性，并说明比较对象怎样选择：
\begin{equation*}
 \text{（1）}\ \sum_{n=1}^{\infty}\frac1{n^2+n+1};
 \qquad
 \text{（2）}\ \sum_{n=1}^{\infty}\frac1{\sqrt{n(n+1)}}.
\end{equation*}
\end{example}

\subsection{比值与根值判别*}

比较对象不必预先知道。若相邻项或第\(n\)次根的变化受到控制，就可以与等比级数建立联系。

\begin{theorem}[比值与根值判别法]
设\(a_n>0\)。若\(a_{n+1}/a_n\to r\)，或\(\sqrt[n]{a_n}\to r\)，则在\(r<1\)时\(\sum a_n\)收敛，在\(r>1\)时发散；\(r=1\)时不能仅据此判断。
\end{theorem}

\begin{proof}
先看比值。当\(r<1\)时，在\(r\)与\(1\)之间取\(q\)，充分靠后的项满足\(a_{n+1}\le qa_n\)，由连乘得\(a_{N+k}\le a_Nq^k\)，故可与等比级数比较。当\(r>1\)时，尾部各项递增且为正，不可能趋于零。

对于根值，\(r<1\)时同样取\(r<q<1\)，充分大的\(n\)满足\(a_n\le q^n\)；\(r>1\)时可取\(1<q<r\)，有\(a_n\ge q^n\)，故通项不趋于零。最后，\(\sum1/n\)与\(\sum1/n^2\)说明\(r=1\)时可能发散，也可能收敛。
\end{proof}

\begin{example}
判断\(\sum_{n=1}^{\infty}n^2/3^n\)、\(\sum_{n=1}^{\infty}2^n/n!\)与\(\sum_{n=1}^{\infty}(n/(2n+1))^n\)的敛散性。比较逐项估计、比值判别与根值判别各自适合的表达式。
\end{example}

\subsection{绝对收敛与交错级数*}

\begin{definition}
若\(\sum|a_n|\)收敛，称\(\sum a_n\){\heiti{}绝对收敛}。若\(\sum a_n\)收敛而\(\sum|a_n|\)发散，称它{\heiti{}条件收敛}。
\end{definition}

\begin{theorem}
绝对收敛的数项级数收敛。
\end{theorem}

\begin{proof}
由\(\left|\sum_{k=n+1}^m a_k\right|\le\sum_{k=n+1}^m|a_k|\)，绝对值级数的Cauchy估计也控制原级数的相应区段，再用式\eqref{eq:analysis-series-cauchy}即可。
\end{proof}

符号交替时，前后项的抵消还可能使不绝对收敛的级数收敛。

\begin{theorem}[交错级数判别法]
若\(a_n\ge0\)单调不增且趋于零，则\(\sum_{n=1}^{\infty}(-1)^{n-1}a_n\)收敛。设其和为\(S\)，前\(n\)项和为\(S_n\)，则
\begin{equation}\label{eq:analysis-alternating-error}
 |S-S_n|\le a_{n+1},
\end{equation}
且\(S\)位于\(S_n\)与\(S_{n+1}\)之间。
\end{theorem}

\begin{proof}
偶数项部分和满足\(S_{2n+2}-S_{2n}=a_{2n+1}-a_{2n+2}\ge0\)，奇数项部分和满足\(S_{2n+3}-S_{2n+1}=-a_{2n+2}+a_{2n+3}\le0\)。又\(S_{2n}\le S_{2n+1}\)，两者之差\(a_{2n+1}\to0\)，故两子列收敛于同一数\(S\)，原部分和数列也收敛于\(S\)。任意相邻两个部分和将\(S\)夹在其中，区间长度为\(a_{n+1}\)，便得余项估计。
\end{proof}

例如\(1-1/2+1/3-1/4+\cdots\)收敛，而其绝对值级数是发散的调和级数，所以它条件收敛。取前\(n\)项近似其和，误差不超过\(1/(n+1)\)，无须先知道和的具体值。判断收敛、求出级数的和、控制近似误差，是相互联系而又不同的任务。

\begin{example}
讨论\(\sum_{n=1}^{\infty}(-1)^{n-1}/n^p\)在实数\(p\)取不同值时的敛散性，并区分绝对收敛与条件收敛。取\(p=2\)，若要求部分和近似的误差不超过\(10^{-4}\)，由上述估计至少应取多少项？
\end{example}

\begin{exercise}[本节习题]
\mbox{}\par\nobreak

（1）求\(\sum_{n=1}^{\infty}1/[(2n-1)(2n+1)]\)的和及前\(n\)项和的误差。

（2）用错位相减求\(\sum_{n=1}^{\infty}(2n-1)/3^n\)的和。

（3）设\(a_n\ge0\)，且\(\sum a_n\)收敛。证明\(\sum a_n^2\)与\(\sum a_n/(1+a_n)\)均收敛。

（4）设\(a_n>0\)，且\(na_n\to c>0\)。判断\(\sum a_n\)的敛散性。若只知\(na_n\to0\)，是否能够确定？

\optionalexerciseitem{5}判断\(\sum_{n=1}^{\infty}n!/n^n\)的敛散性。再举两个比值极限都等于\(1\)、但敛散性不同的正项级数。

\optionalexerciseitem{6}设\(a_n\ge0\)单调不增且趋于零。证明\(\sum_{n=1}^{\infty}(-1)^{n-1}a_n\)的第\(n\)项余项与第\(n+1\)项同号，余项为零时除外。
\end{exercise}

\newpage
\section{两个重要的常数}
\label{sec:analysis-constants}

\subsection{自然对数的底\texorpdfstring{\(e\)}{e}}
\(e\)是通过数列\(\left\{ e_{n} \right\}\)的极限引进的一个常数，\(e \approx 2.71828\)，其中\(e_{n} = {\left( 1 + \frac{1}{n} \right)}^{n}\).

\begin{example}
已知正数\(a\)，将其分成若干部分，若要使它们的乘积达到最大，应该怎样分？
\end{example}

\subsection{关于数\texorpdfstring{\(e\)}{e}的几个基本命题}
设\(e_{n} = {\left( 1 + \frac{1}{n} \right)}^{n} , n \in \mathbb{N}^{*}\)，证明：\(\left\{ e_{n} \right\}\)严格单调递增且收敛.

（1）证明：\(e = \sum _{n = 0}^{\infty} \frac{1}{n !}\).

（2）记\(\varepsilon_{n} = e - \sum _{k = 0}^{n} \frac{1}{k !}\)，证明：\(\lim_{n \to \infty} \varepsilon_{n} \left( n + 1 \right) ! = 1\).

（3）对于上述\(\varepsilon_{n}\)有\(\frac{1}{\left( n + 1 \right) !} < \varepsilon_{n} < \frac{1}{n ! n}\).

（4）\(e\)是无理数.

（5）\(\lim_{n \to \infty} \frac{n}{\sqrt[n]{n !}} = e\).

\subsection{Euler常数\texorpdfstring{\(\gamma\)}{gamma}}
\label{subsec:analysis-euler-constant}
证明：数列\(\left\{ \gamma_{n} \right\}\)收敛，其中\(\gamma_{n} = \sum _{i = 1}^{n} \frac{1}{i} - \ln n , n \in \mathbb{N}^{*}\)。其极限称为{\heiti{}Euler常数}，记为\(\gamma\)。

计算极限：
\begin{equation}
    \lim_{n \to \infty} \left( \frac{1}{n + 1} + \frac{1}{n + 2} + \cdots + \frac{1}{2 n} \right)
\end{equation}

{\kaiti{}思考\quad Euler常数\(\gamma\)是否是无理数？}

\begin{example}
    计算以下极限：

    （1）\(\lim_{n \to \infty} {\left( 1 - \frac{1}{n} \right)}^{n}\),
    \(\lim_{n \to \infty} {\left( 1 + \frac{1}{2 n} \right)}^{n}\)

    （2）\(\lim_{n \to \infty} {\left( 1 + \frac{1}{n} \right)}^{n^{2}}\)

    （3）\(\lim_{n \to \infty} {\left( 1 + \frac{1}{n} + \frac{1}{n^{2}}
        \right)}^{n}\)

    （4）\(\lim_{n \to \infty} \prod _{i = 1}^{n} \left( 1 +
    \frac{i}{n^{2}} \right)\)

    （5）\(\lim_{n \to \infty} \frac{n ! 2^{n}}{n^{n}}\)

    （6）\(\lim_{n \to \infty} {\left( 1 +
        \frac{1}{p_{n}}\right)}^{p_{n}}\)，其中\(\left\{ p_{n} \right\}\)是正数列
\end{example}


\section{迭代产生的数列}
\label{sec:analysis-iterations}
所谓由迭代产生的数列，是指在给出数列的第一项\(a_{1}\)后，用递推公式\(a_{n + 1} = f \left( a_{n} \right) , n \in \mathbb{N}^{*}\)通过迭代生成的数列。迭代是近似计算方法的重要实现过程，在数学和其他领域经常出现，有很强的理论和实用价值。

\subsection{几个基础的例子}\label{subsec:iteration-examples}
\paragraph{1°逐项开方的迭代} \mbox{}

设\(a_{1} = \sqrt{2} , a_{n + 1} = \sqrt{2 + a_{n}} , n \in \mathbb{N}^{*}\).讨论\(\left\{ a_{n} \right\}\)的敛散性。若收敛，求出其极限。

可以归纳地证明这个数列是严格单调递增的，且以2为上界。

显然，由递推公式和初始值知，该数列的每一项均为正。由于
\begin{equation}
    a_{1} = \sqrt{2} < \sqrt{2 + \sqrt{2}} = a_{2}
\end{equation}
以及当\(a_{n - 1} < a_{n}\)时，有
\begin{equation}
    a_{n} = \sqrt{2 + a_{n - 1}} < \sqrt{2 + a_{n}} = a_{n + 1}
\end{equation}
严格成立，可见数列是严格单调递增的。又由\(a_{1} < 2\)以及当\(a_{n} < 2\)时，有
\begin{equation}
    a_{n + 1} = \sqrt{2 + a_{n}} < \sqrt{2 + 2} = 2
\end{equation}
可知\(\forall n \in \mathbb{N}^{*} , a_{n} < 2\)成立。即数列以2为上界。因此该数列收敛，记极限为\(a\)。在递推式\(a_{n+1}=\sqrt{2+a_n}\)两边取极限，得到
\begin{equation}
    a = \sqrt{2 + a} , 0 < a \le 2
\end{equation}
解得\(a = 2\)，即为所求之\(\left\{ a_{n} \right\}\)极限。

\paragraph{2°倒数迭代} \mbox{}

设数列\(\left\{ b_{n} \right\}\)满足\(b_{1} = 1 , b_{n + 1} = 1 + \frac{1}{b_{n}} \left( n \in \mathbb{N}^{*} \right)\)。讨论\(\left\{ b_{n} \right\}\)的敛散性，若收敛，求出其极限。

显然，由递推公式和初始值知，该数列的每一项均为正。先假定该数列收敛到\(b\)，则仿照上例有\(b = 1 + \frac{1}{b} \to b = \frac{1 + \sqrt{5}}{2}\)。由递推公式得\(b_{n + 2} = 2 - \frac{1}{\left( b_{n} + 1 \right)}\)，故有
\begin{equation}
    b_{n + 2} - b_{n} = \frac{b_{n} - b_{n - 2}}{\left( b_{n} + 1 \right) \left( b_{n - 2} + 1 \right)}
\end{equation}
容易看出\(b_{n + 2} - b_{n}\)与\(b_{n} - b_{n - 2}\)同号，故\(\left\{ b_{n} \right\}\)的奇子列和偶子列各自单调。
不难计算出数列的前几项\(b_{1} = 1 , b_{2} = 2 , b_{3} = \frac{3}{2} , b_{4} = \frac{5}{3}\)，可见\(\left\{ b_{n} \right\}\)的奇子列严格单调递增，偶子列严格单调递减，由归纳法易证
\begin{equation}
    b_{1} < b_{3} < \cdots < b < \cdots < b_{4} < b_{2}
\end{equation}
因此\(\left\{ b_{n} \right\}\)的奇子列和偶子列分别单调有界，均收敛。由\(b_{n + 2} = 2 - \frac{1}{b_{n} + 1}\)，在两个子列上分别取极限，可知它们的极限均满足\(l = 2 - \frac{1}{l + 1}\)。又两个极限均为正，故均等于\(b = \frac{1 + \sqrt{5}}{2}\)。故\(\left\{ b_{n} \right\}\)收敛到\(b\)，为\(\frac{1 + \sqrt{5}}{2}\)。

\paragraph{Fibonacci数列} \mbox{}

本题与Fibonacci数列有关。所谓Fibonacci数列，即是由
\begin{equation}
    F_{1} = 1 , F_{2} = 1 , F_{n + 2} = F_{n + 1} + F_{n} , n \in \mathbb{N}^{*}
\end{equation}
迭代而成的数列。如计算相邻两项的比值（或称为增长率）的极限，即\(\lim_{n \to \infty} \frac{F_{n + 1}}{F_{n}}\)，则
\begin{equation}
    \frac{F_{n + 1}}{F_{n}} = \frac{F_{n} + F_{n - 1}}{F_{n}} = 1 + \frac{1}{\frac{F_{n}}{F_{n - 1}}}
\end{equation}

定义\(b_{n} = \frac{F_{n + 1}}{F_{n}}\)，则显然有\(b_{1} = 1 , b_{n + 1} = 1 + \frac{1}{b_{n}}\).对应本例题可知其极限为\(\frac{\sqrt{5} + 1}{2}\)。

\enlargethispage{\baselineskip}
\subsection{单调性与几何方法}
关于迭代生成数列的第一个规律可以用如下命题表示。

\begin{theorem}
（第一律）设数列\(\left\{ x_{n} \right\}\)满足递推公式\(x_{n + 1} = f \left( x_{n} \right) , n \in \mathbb{N}^{*}\)。若有\(\lim_{n \to \infty} x_{n} = \xi\)，同时又成立
\begin{equation}
    \lim_{n \to \infty} f \left( x_{n} \right) = f \left( \xi \right)
\end{equation}
则极限\(\xi\)一定是方程\(f \left( x \right) = x\)的根（这时称\(\xi\)为函数\(f\)的{\heiti{}不动点}）。
\end{theorem}

该命题的证明是平凡的，只需在递推公式的两边取\(n \to \infty\)即可。后面我们会知道该命题的条件可替换为\(f\)在点\(\xi\)处连续等条件。

这个命题的用处是明显的。在我们不知道数列是否收敛的情况下，可以先求解方程\(f \left( x \right) = x\)。求出方程的根对判定原数列是否收敛是有帮助的。例如，如果\(f\)在\(\mathbb R\)上连续，而方程\(f \left( x \right) = x\)在实数范围内无根，则无需再做任何研究就可以断言：这个由迭代产生的数列一定发散。在上述\ref{subsec:iteration-examples}节中的第二个例子中，一开始求出\(b\)，到最后才证明它是极限，我们已经看到在证明中\(b\)所起的作用。
\newpage
\begin{theorem}
（第二律）设\(\left\{ x_{n} \right\}\)满足关系\(x_{n + 1} = f \left( x_{n} \right) , n \in \mathbb{N}^{*}\)，其中的函数\(f\)在区间\(I\)上单调，同时数列\(\left\{ x_{n} \right\}\)的每一项都在区间\(I\)中，则只有两种可能：

（1）当\(f\)为单调增加时，\(\left\{ x_{n} \right\}\)为单调数列；

（2）当\(f\)为单调减少时，\(\left\{ x_{n} \right\}\)的两个子列\(\left\{ x_{2 k - 1} \right\}\)和\(\left\{ x_{2 k} \right\}\)分别为单调数列，且具有相反的单调性。
\end{theorem}

对于第一种可能，设\(f \left( x \right)\)在区间\(I\)上单调增加。根据条件，有\(x_{n} \in I , n \in \mathbb{N}^{*}\)。观察数列的前两项，如果\(x_{1} \le x_{2} = f \left( x_{1} \right)\)，则有\(x_{2} = f \left( x_{1} \right) \le f \left( x_{2} \right) = x_{3}\)。由归纳法，知数列\(\left\{ x_{n} \right\}\)单调增加。类似地，当\(x_{1} \ge x_{2}\)时，数列\(\left\{ x_{n} \right\}\)单调减少。

对于第二种可能，设\(f \left( x \right)\)在区间\(I\)上单调减少。容易知道，复合函数\(f \left( f \left( x \right) \right)\)是单调增加的。严格地说，\(\forall a , b \in I , a < b , f \left( a \right) , f \left( b \right) \in I \Rightarrow f \left( f \left( a \right) \right) \le f \left( f \left( b \right) \right)\)。

观察\(x_{1}\)与\(x_{3}\)。如果\(x_{1} = x_{3}\)，则子列\(\left\{ x_{2 k - 1} \right\}\)为常数列。如果\(x_{1} < x_{3}\)，从\(f\)单调减少易推知\(x_{2} = f \left( x_{1} \right) \ge f \left( x_{3} \right) = x_{4}\)，然后推知\(x_{3} \le x_{5}\)。于是用数学归纳法易证此时\(\{ x_{2 k - 1} \}\)单调递增，由于\(f\)单调递减且\(\left\{ x_{2 k} \right\} = \left\{ f \left( x_{2 k - 1} \right) \right\}\)单调递减。对于\(x_{1} > x_{3}\)的情况，讨论完全类似，从略。

\begin{theorem}\label{thm:iter-fixed}
设\(a\)是\(f \left( x \right)\)的不动点，函数\(f\)在\(a\)的邻域\(\left( a - r , a + r \right)\)中连续且严格单调递增，且在\(\left( a - r , a \right)\)上有\(f \left( x \right) > x\)，而区间\(\left( a , a + r \right)\)上有\(f \left( x \right) < x\)，那么由迭代产生的数列\(x_{n + 1} = f \left( x_{n} \right)\)只要第一项在\(\left( a - r , a + r \right)\)内，且不等于\(a\)，则以后就不会越出这个区间，而且是以\(a\)为极限的严格单调数列。
\end{theorem}

定理\ref{thm:iter-fixed}的证明留给读者自行证明。

\begin{example}
    完成以下计算或证明：
    
    （1）设\(a > 0 , x_{1} = \sqrt{a} , x_{n + 1} = \sqrt{a + x_{n}} , n \in \mathbb{N}^{*}\)，求\(\lim_{n \to \infty} x_{n}\).

    （2）设\(a > 0 , x_{1} = \sqrt{a} , x_{n + 1} = \sqrt{a x_{n}} , n \in \mathbb{N}^{*}\)，求\(\lim_{n \to \infty} x_{n}\).

    （3）设\(x_{1} = b , x_{n + 1} = \frac{1}{2} \left( x_{n}^{2} + 1 \right)\)，问：\(b\)取何值时数列\(\left\{ x_{n} \right\}\)收敛，并求其极限。

    （4）设\(\left\{ x_{n} \right\}\)为正数列，且满足\(x_{n + 1} + \frac{1}{x_{n}} < 2\)，证明：\(\left\{ x_{n} \right\}\)收敛，并求其极限。

    （5）设\(A > 0 , x_{1} > 0 , x_{n + 1} = \frac{x_{n} \left( x_{n}^{2} + 3 A \right)}{3 x_{n}^{2} + A} , n \in \mathbb{N}^{*}\)，证明：\(\left\{ x_{n} \right\}\)收敛于\(\sqrt{A}\)。
\end{example}



\newpage
\section{综合习题}
\label{sec:analysis-comprehensive-exercises}

以下习题覆盖数列的表示、求和、极限、数项级数与迭代。求通项时应说明依据；研究极限时，应区分极限的候选值与收敛性的证明。

\begin{exercise}
\label{exs:analysis-comprehensive}
\mbox{}\par\nobreak

（1）设\(a_1=1,a_{n+1}-a_n=2n\)。

\quad \exercisesubitem{1}用累加法求通项，并用数学归纳法核对。

\quad \exercisesubitem{2}求前\(n\)项和\(S_n\)，再由\(S_n-S_{n-1}\)恢复\(a_n\)。

（2）等差数列的前\(n\)项和为\(S_n\)，且\(S_3=12,S_6=42\)。求通项，并求
\begin{equation*}
 T_n=\sum_{k=1}^n\frac1{a_ka_{k+1}}
\end{equation*}
及\(\lim_{n\to\infty}T_n\)。

（3）等比数列的首项为\(3\)，公比为\(-1/2\)。求通项、前\(n\)项和及二者的极限。原数列与部分和数列是否单调？它们的奇、偶子列分别怎样变化？

（4）设\(S_n=\sum_{k=1}^n(2k-1)3^{k-1}\)。

\quad \exercisesubitem{1}用错位相减求\(S_n\)。

\quad \exercisesubitem{2}分别求\(S_n/3^n\)、\(S_n/(n3^n)\)的极限，比较两种尺度下的结果。

（5）设
\begin{equation*}
 S_n=\sum_{k=1}^n\frac1{\sqrt{k+1}+\sqrt k}.
\end{equation*}
求\(S_n\)，并分别求\(\lim_{n\to\infty}S_n/n\)、\(\lim_{n\to\infty}S_n/\sqrt n\)。说明同一个部分和数列在不同尺度下的极限有何不同。

（6）\exercisesubitem{1}只用定义证明\((3n-1)/(n+2)\to3\)，并对给定的\(\varepsilon>0\)写出一个可用的\(N\)。

\quad \exercisesubitem{2}设\(a_n\to a,b_n\to b\)，且从某项起\(a_n\le b_n+1/n\)，证明\(a\le b\)。

（7）设实数列满足\(|a_{n+1}-a_n|\le2^{-n}\)。用Cauchy准则证明它收敛，并给出极限\(a\)与\(a_n\)之间的误差上界。若条件换为\(|a_{n+1}-a_n|\le1/n\)，是否仍必收敛？说明理由。

（8）若数列\(\{a_n\}\)的奇子列与偶子列都收敛于\(a\)，证明原数列也收敛于\(a\)。

（9）从\([a_0,b_0]=[1,2]\)开始，每次取\(m_n=(a_n+b_n)/2\)。若\(m_n^3\le2\)，保留\([m_n,b_n]\)；否则保留\([a_n,m_n]\)，记为\([a_{n+1},b_{n+1}]\)。

\quad \exercisesubitem{1}证明这些区间构成长度趋于零的区间套，公共点\(\xi\)满足\(\xi^3=2\)。

\quad \exercisesubitem{2}用进行十次二分后的区间中点近似\(\xi\)，给出误差上界。这里的收敛性为何不需要先求出\(\xi\)？

（10）利用适当的极限工具，计算或证明：
\begin{equation*}
 \lim_{n\to\infty}\frac{\sqrt1+\sqrt2+\cdots+\sqrt n}{n^{3/2}}=\frac23;
 \qquad
 x_n\to l\ \Longrightarrow\
 \frac{\sum_{k=1}^n kx_k}{n^2}\to\frac l2.
\end{equation*}
说明相邻增量的比较怎样确定上述极限。

（11）记\(E_n=\sum_{k=0}^n1/k!\)。

\quad \exercisesubitem{1}由\(0<e-E_n<1/(n\,n!)\)，求使这个误差上界小于\(10^{-6}\)的最小正整数\(n\)。

\quad \exercisesubitem{2}利用本章关于余项的结论，证明\(n\,n!(e-E_n)\to1\)。误差上界与实际误差的比值趋于多少？

（12）设\(A>0,x_1>0\)，并令
\begin{equation*}
 x_{n+1}=\frac12\left(x_n+\frac A{x_n}\right).
\end{equation*}
证明从第二项起\(x_n\ge\sqrt A\)，且数列单调不增，进而证明\(x_n\to\sqrt A\)。再证明误差关系
\begin{equation*}
 x_{n+1}-\sqrt A=\frac{(x_n-\sqrt A)^2}{2x_n},
\end{equation*}
说明它为何适合用于平方根的近似计算。

（13）计算下列极限，并说明怎样消去式中的主要部分：
\begin{equation*}
 \exercisesubitem{1}\ \lim_{n\to\infty}\left(\sqrt{n^2+n}-n\right);
 \qquad
 \exercisesubitem{2}\ \lim_{n\to\infty}n\left(\sqrt{n^2+n}-n-\frac12\right).
\end{equation*}

（14）设\(a_n\to l\)。证明
\begin{equation*}
 \lim_{n\to\infty}\frac{a_1+2a_2+\cdots+2^{n-1}a_n}{2^n}=l.
\end{equation*}
若将条件改为左侧数列收敛于\(l\)，是否能反推出\(a_n\to l\)？与Cauchy命题的逆命题比较。

（15）设\(\alpha>-1\)。求\(\sum_{n=1}^{\infty}1/[(n+\alpha)(n+\alpha+1)]\)的和，并求前\(N\)项和的误差。对给定的\(\varepsilon>0\)，写出保证误差小于\(\varepsilon\)的\(N\)应满足的条件。

（16）设\(a_n>0\)单调不增，且\(\sum_{n=1}^{\infty}a_n\)收敛。

\quad \exercisesubitem{1}通过比较\(na_{2n}\)与\(\sum_{k=n+1}^{2n}a_k\)，证明\(na_n\to0\)。

\quad \exercisesubitem{2}若不要求单调，结论是否仍成立？试用只在少数下标处较大的数列说明。

（17）设\(x_1>0\)，且\(x_{n+1}=x_n/(1+x_n)\)。

\quad \exercisesubitem{1}考察\(1/x_n\)，求\(x_n\)的通项与\(\lim_{n\to\infty}nx_n\)。

\quad \exercisesubitem{2}证明\(x_n-x_{n+1}=x_nx_{n+1}\)，并求级数\(\sum_{n=1}^{\infty}x_nx_{n+1}\)的和与前\(N\)项和的误差。

（18）设\(\{a_n\}\)有界，且它的每个收敛子列都以\(a\)为极限。证明\(a_n\to a\)。若原数列不收敛于\(a\)，试先找出始终与\(a\)相隔一定距离的子列，再使用有界数列的收敛子列定理。

（19）设\(S_n=\sum_{k=1}^n a_k\)有界。

\quad \exercisesubitem{1}证明\((a_1+\cdots+a_n)/n\to0\)。这是否说明\(a_n\to0\)，或说明\(\sum a_n\)收敛？分别说明理由。

\quad \exercisesubitem{2}若另有\(a_n\ge0\)，结论有什么变化？比较有界性与单调性在这里的作用。

（20）设闭区间\([u_{n+1},v_{n+1}]\subseteq[u_n,v_n]\)，且\(v_n-u_n\to d\ge0\)。

\quad \exercisesubitem{1}证明两端点数列分别收敛，记极限为\(u,v\)，并证明全部区间的公共部分恰为\([u,v]\)，其长度为\(d\)。

\quad \exercisesubitem{2}若改为开区间，即使长度趋于零，是否仍一定有公共点？举例说明。

（21）记\(H_n=\sum_{k=1}^n1/k\)，已知\(\gamma_n=H_n-\ln n\)收敛。对固定整数\(m\ge2\)，求\(\lim_{n\to\infty}(H_{mn}-H_n)\)。为什么不能把含有\(n\)项的这个和，直接拆成逐项极限之和？
\end{exercise}

\begin{exercise}[选做]
\label{exs:analysis-comprehensive-optional}
\mbox{}\par\nobreak

（1）判断\(\sum_{n=1}^{\infty}(-1)^{n-1}(\sqrt{n+1}-\sqrt n)\)是否收敛、是否绝对收敛，并给出前\(N\)项和近似级数和的误差上界。若要求误差不超过\(10^{-2}\)，应怎样选择\(N\)？
\end{exercise}

